\subsection{Dataset and Data Collection}
\label{subsec:evaluation_dataset}

The main training and synthetic evaluation benchmark was generated using a
controlled Unreal Engine simulation that was modified and automated for this
project. The simulator rendered monocular \ac{RGB} images of a single target
drone at known metric distances while varying the environmental conditions.
Each rendered image was paired with one YOLO-format annotation file containing
the normalized bounding box of the target drone.

The synthetic metadata were encoded directly in the filename using a fixed
template:
\begin{center}
\small
\texttt{<frame>\_depth\_<distance>\_<weather>\_<time>.png}.
\end{center}
During preprocessing, the pipeline parsed this filename pattern to recover the
ground-truth distance, weather condition, and time of day for each image. The
paired YOLO label file provided the target-drone bounding box used for
geometric feature extraction.

The final unified synthetic benchmark contains 15,064 images. It covers 12
discrete ground-truth distances:
\[
20, 30, 40, 50, 60, 70, 80, 90, 100, 115, 125, 150 \text{ m}.
\]
For analysis, these distances were grouped into three operating regimes:
near range ($z \leq 60$~m), mid range ($60 < z \leq 100$~m), and far range
($z > 100$~m). The synthetic benchmark includes two weather conditions,
\texttt{clear\_sky} and \texttt{light\_rain}, and two time-of-day settings,
\texttt{10AM} and \texttt{8PM}. The final experiments were conducted on still
images; video sequences were not used for training or evaluation.

Ground-truth distance in the synthetic benchmark was obtained from the
controlled rendering configuration. Because the simulator explicitly placed the
target drone at each specified depth, the filename and folder structure provide
metric distance labels without requiring manual measurement. The bounding-box
annotations were used only to localize the drone in the image and to derive
image-space geometry features such as box size, aspect ratio, and image
position.

In addition to the synthetic benchmark, AirDepth was evaluated on a prepared
real-drone external dataset derived from the Nenrus drone-distance dataset
\cite{nenrus_drone_distance_dataset}. This real-image set contains 489 images
from two subsets, Kongsberg and Vestfold, with 203 and 286 samples,
respectively. The images have resolution $1920 \times 1080$ pixels and include
ground-truth distance labels from 4~m to 75~m. These images were not used to
train the synthetic AirDepth model. They were used only to evaluate
synthetic-to-real transfer and to fit the lightweight real-world calibration
stage described in Section~\ref{subsec:robustness_analysis}.

\begin{table}[t]
\caption{Evaluation datasets.}
\label{tab:evaluation_datasets}
\centering
\small
\setlength{\tabcolsep}{4pt}
\renewcommand{\arraystretch}{1.15}
\begin{tabular}{lccp{2.9cm}}
\hline
\textbf{Dataset} &
\textbf{Images} &
\textbf{Range} &
\textbf{Conditions / purpose} \\
\hline
Synthetic benchmark &
15,064 &
20--150~m &
Controlled Unreal Engine scenes with clear/rain weather and 10~AM/8~PM lighting \\
Nenrus real-drone &
489 &
4--75~m &
Real outdoor images used for external transfer evaluation and calibration \\
\hline
\end{tabular}
\end{table}

\section{Dataset Documentation}
\label{sec:datasheet}

\noindent This appendix documents the synthetic \airdepth\ drone-distance
dataset introduced in Section~\ref{subsec:evaluation_dataset}, following the
\emph{Datasheets for Datasets} framework of Gebru \etal~\cite{gebru2021datasheets}.

% -------------------------------------------------------------------------
\subsection{Motivation}

\begin{description}
    \item[For what purpose was the dataset created?] \hfill \\{}
    The dataset was created to support supervised monocular drone-distance
    estimation under controlled and repeatable conditions. It addresses the
    lack of readily available UAV image datasets with metric ground-truth
    distance labels, synchronized target-drone bounding boxes, and explicit
    environmental metadata. The controlled simulation also enables systematic
    evaluation across distance, weather, and time-of-day regimes.
\end{description}

% -------------------------------------------------------------------------
\subsection{Composition}

\begin{description}
    \item[What do the instances represent?] \hfill \\{}
    Each instance is a single rendered monocular \ac{RGB} frame of one target
    drone at a known metric distance, paired with a YOLO-format bounding-box
    annotation.

    \item[How many instances are there in total?] \hfill \\{}
    15,064 source images, partitioned into 12,804 development samples and
    2,260 held-out test samples (85/15 split).

    \item[Does the dataset contain all possible instances or a sample?] \hfill \\{}
    The dataset is a finite synthetic benchmark rendered for this project. It
    should be treated as a sampled coverage of the selected simulation
    conditions rather than an exhaustive enumeration of all possible drone
    poses, backgrounds, weather states, lighting conditions, or distances.

    \item[What data does each instance consist of?] \hfill \\{}
    An \ac{RGB} image; a YOLO-format normalized bounding box for the target
    drone; and ground-truth metadata (distance, weather condition, time of
    day) encoded in the filename using the template
    \texttt{<frame>\_depth\_<distance>\_<weather>\_<time>.png}.

    \item[Is there a label or target associated with each instance?] \hfill \\{}
    Yes: the ground-truth metric distance $z$ (one of 12 discrete values:
    20, 30, 40, 50, 60, 70, 80, 90, 100, 115, 125, 150~m).

    \item[Is any information missing from individual instances?] \hfill \\{}
    No missing fields are expected for the curated benchmark: each retained
    image has a paired bounding-box label and distance, weather, and
    time-of-day metadata recoverable from the filename. Frames that do not
    follow the expected filename or label format are excluded from the feature
    extraction tables.

    \item[Are relationships between instances made explicit?] \hfill \\{}
    Related instances are made explicit through the structured file paths and
    filename metadata. Images are grouped by distance, weather, and time of day,
    and downstream feature tables preserve a source-image identifier so that
    augmented or noisy-bounding-box variants derived from the same source image
    remain traceable.

    \item[Are there recommended data splits?] \hfill \\{}
    Yes. The recommended split is the fixed 85/15 development/test split used
    in the experiments, with 12,804 development images and 2,260 held-out test
    images. For augmented experiments, all rows derived from the same source
    image are kept in the same split group to avoid leakage between train and
    test data.

    \item[Are there any errors, sources of noise, or redundancies?] \hfill \\{}
    The distance labels are exact by construction, but the dataset can still
    contain simulation-domain bias, visually similar frames within the same
    condition group, quantization from the 12 discrete distance values, and
    possible annotation or parsing errors in malformed image/label pairs.
    Additional detector-like noise was introduced only in specific robustness
    studies by perturbing bounding boxes; it is not part of the base rendered
    image set.

    \item[Is the dataset self-contained, or does it rely on external resources?] \hfill \\{}
    The released synthetic image/label benchmark is self-contained. Recreating
    the images from scratch would require the original Unreal Engine simulation
    assets and rendering scripts.

    \item[Does the dataset contain data that might be considered confidential?] \hfill \\{}
    No.

    \item[Does the dataset contain data that might be offensive, insulting, threatening, or anxiety-inducing?] \hfill \\{}
    No.
\end{description}

% -------------------------------------------------------------------------
\subsection{Collection Process}

\begin{description}
    \item[How was the data associated with each instance acquired?] \hfill \\{}
    Rendered directly via a controlled Unreal Engine simulation rather than
    measured or inferred from real-world sensors. The ground-truth distance is
    exact by construction because the simulator places the target drone at the
    specified metric depth.

    \item[What mechanisms or procedures were used to collect the data?] \hfill \\{}
    The collection process used an automated Unreal Engine rendering setup. For
    each configured scene condition, the simulator positioned a single target
    drone at one of the predefined metric distances and rendered a monocular
    \ac{RGB} frame. The same process exported the corresponding YOLO-format
    target-drone bounding box and embedded the distance, weather, and
    time-of-day metadata in the image filename.

    \item[If the dataset is a sample, what was the sampling strategy?] \hfill \\{}
    The benchmark samples the Cartesian set of selected distances, weather
    conditions, and time-of-day settings, with multiple rendered frames per
    condition. The selected distances emphasize the operational range relevant
    to the project, from 20~m to 150~m.

    \item[Who was involved in the data collection process?] \hfill \\{}
    The dataset was generated and curated by the project authors as part of the
    AirDepth final project workflow.

    \item[Over what timeframe was the data collected?] \hfill \\{}
    The synthetic data were rendered during the development of this project,
    before the final training and evaluation studies were run.

    \item[Were any ethical review processes conducted?] \hfill \\{}
    Not applicable; the dataset is fully synthetic and contains no human
    subjects or personal data.
\end{description}

% -------------------------------------------------------------------------
\subsection{Preprocessing, Cleaning, and Labeling}

\begin{description}
    \item[Was any preprocessing/cleaning/labeling of the data done?] \hfill \\{}
    Yes. The pipeline parsed each filename to extract distance, weather, and
    time-of-day metadata; matched each image with its YOLO-format label file;
    converted normalized bounding boxes into image-space geometry features; and
    removed rows that could not be parsed or matched consistently. Later
    experiments also derived relative-depth features inside the target
    bounding box and expanded context windows.

    \item[Was the ``raw'' data saved in addition to the preprocessed data?] \hfill \\{}
    Yes. The raw rendered images and YOLO label files are preserved separately
    from derived feature tables, split files, model reports, and evaluation
    artifacts.

    \item[Is the software used to preprocess/clean/label the data available?] \hfill \\{}
    The repository contains the Python code used for feature extraction,
    dataset loading, split construction, calibration, and evaluation. The
    original Unreal Engine rendering project and simulation assets are separate
    from the derived experimental code.
\end{description}

% -------------------------------------------------------------------------
\subsection{Uses}

\begin{description}
    \item[Has the dataset been used for any tasks already?] \hfill \\{}
    Yes, for training and evaluating the \airdepth\ monocular
    distance-estimation framework presented in this paper (see
    Section~\ref{sec:evaluation}).

    \item[Is there a repository that links to any papers or systems that use the dataset?] \hfill \\{}
    The dataset is used by the AirDepth project repository associated with this
    paper. No independent external systems using the synthetic benchmark are
    known at the time of writing.

    \item[What other tasks could the dataset be used for?] \hfill \\{}
    The dataset may also support drone detection, monocular distance-estimation
    baselines, bounding-box geometry studies, synthetic-to-real adaptation
    experiments, and robustness analysis under controlled weather and lighting
    changes.

    \item[Is there anything about the composition or collection that might impact future uses?] \hfill \\{}
    Yes. The dataset is synthetic, contains a single target drone per image,
    uses 12 discrete distance values, and covers only two weather conditions
    and two time-of-day settings. Results obtained on it may not generalize
    directly to other drone shapes, camera intrinsics, cluttered real
    backgrounds, occlusions, motion blur, adverse weather, or safety-critical
    operational settings without additional real-world validation.

    \item[Are there tasks for which the dataset should not be used?] \hfill \\{}
    The dataset should not be used as the sole validation source for
    safety-critical drone-distance estimation or collision-avoidance
    deployment. It should also not be treated as a real-world UAV benchmark;
    the external Nenrus evaluation and calibration results are needed to assess
    synthetic-to-real transfer.
\end{description}
